% ============================================================================
%  presentation/script-short.tex --- the 30-seconds-a-slide script, EN | FR
%
%  The same talk as script.tex, cut to fit. Thirty-six content slides at half a
%  minute each, plus eight section signposts of about five seconds: eighteen and
%  a half minutes, against roughly thirty for the full script.
%
%  WHAT "30 SECONDS" MEANS HERE. At an unhurried 150 words a minute, thirty
%  seconds is about seventy-five words. The cells below average SIXTY, and none
%  exceeds seventy-five -- roughly twenty-five seconds of speech, which leaves
%  five seconds a slide for the pause before you start, for pointing at the
%  figure, and for the breath after the sentence that matters. A cell written to
%  the full seventy-five is a cell read at a rush. Add nothing here without
%  taking something out.
%
%  WHAT WAS CUT, AND WHAT WAS NOT. The caveats stay. Every one of them. What
%  goes is elaboration: the second example, the reason behind the reason, the
%  aside that makes a point land harder rather than land at all. A talk can be
%  short or long; it cannot be short and overclaiming, because the sentence that
%  gets dropped under time pressure is always the qualifying one.
%
%  Layout and shorthands come from scriptstyle.tex, shared with script.tex.
%  Slide titles are verbatim from main.tex and enforced by check_script.py.
%
%  Build:  make -C presentation script-short   ->  script-short.pdf
% ============================================================================
\documentclass[10pt,a4paper]{article}
% ============================================================================
%  presentation/scriptstyle.tex --- shared preamble for the speaker scripts
%
%  script.tex (the full one) and script-short.tex (30 seconds a slide) are the
%  same document at two lengths. They looked the same by coincidence until this
%  file existed; now they look the same by construction, and a change to the
%  two-column layout cannot land in one and not the other.
%
%  \scriptname{...} sets the right-hand running head, which is the only thing
%  the two documents disagree about.
% ============================================================================
% Both \scriptname and its default live inside \makeatletter: defining the
% setter outside it made LaTeX read \@scriptname as \@ followed by the letters
% "scriptname", so the setter silently did nothing and both documents kept the
% default running head.
\makeatletter
\newcommand{\scriptname}[1]{\gdef\@scriptname{#1}}
\gdef\@scriptname{EN spoken \; $\cdot$ \; FR pour comprendre}
\makeatother

\usepackage{iftex}
\ifLuaTeX
  \usepackage{fontspec}
  \IfFontExistsTF{Lato}{\setsansfont{Lato}}{}
\else
  \usepackage[utf8]{inputenc}
  \usepackage[T1]{fontenc}
  \usepackage{lmodern}
\fi
\usepackage[english,french]{babel}
\usepackage[a4paper,margin=15mm,top=17mm,bottom=17mm]{geometry}
\usepackage{xcolor}
\usepackage{array}
\usepackage{tabularx}
\usepackage{booktabs}
\usepackage{microtype}
\usepackage{enumitem}
\usepackage{fancyhdr}
\usepackage[hidelinks]{hyperref}

\definecolor{ink}   {HTML}{1A1A1A}
\definecolor{accent}{HTML}{1F5C8B}
\definecolor{good}  {HTML}{2E7D32}
\definecolor{bad}   {HTML}{C62828}
\definecolor{warn}  {HTML}{EF6C00}
\definecolor{muted} {HTML}{6E6E6E}
\definecolor{rule}  {HTML}{D8D8D4}
\definecolor{wash}  {HTML}{F4F6F8}

\renewcommand{\familydefault}{\sfdefault}
\color{ink}
\setlength{\parindent}{0pt}
\setlength{\parskip}{0pt}

\pagestyle{fancy}
\fancyhf{}
\renewcommand{\headrulewidth}{0pt}
\fancyfoot[C]{\footnotesize\color{muted}\thepage}
\fancyhead[L]{\footnotesize\color{muted}Noise Modelling Hanoi --- speaker script}
\makeatletter\fancyhead[R]{\footnotesize\color{muted}\@scriptname}\makeatother

% --------------------------------------------------------------- structure
% \slide{number}{short title}{minute mark}
% \slide[<on-screen fraction>]{<page>}{<title>}{<minute mark>}
%
% TWO NUMBERS, because the deck shows one and the file has the other. \#1 is the
% PDF page -- what you get by pressing right that many times from the start. The
% optional argument is the FRACTION PRINTED ON THE SLIDE ITSELF, which is lower,
% because metropolis does not count section pages or the appendix. Presenters
% navigate by what they can see, so the visible number is given first and in
% bold; the page number follows it in grey for anyone opening the PDF.
\newcommand{\slide}[4][]{%
  \par\vspace{5mm}\noindent
  \colorbox{accent}{\parbox{\dimexpr\textwidth-2\fboxsep}{%
    \color{white}\vspace{.6mm}
    \if\relax\detokenize{#1}\relax
      \textbf{#2}%
    \else
      \textbf{#1}\,{\footnotesize\color{white!70}(p.\,#2)}%
    \fi
    \hspace{3mm}#3\hfill\small #4\vspace{.6mm}}}
  \par\vspace{1.5mm}}

% \section-level marker for an act of the talk
\newcommand{\act}[2]{%
  \par\vspace{7mm}\noindent
  {\color{accent}\rule{\textwidth}{.9pt}}\par\vspace{1.5mm}
  {\large\bfseries\color{accent}#1}\hfill{\color{muted}#2}
  \par\vspace{.5mm}{\color{rule}\rule{\textwidth}{.4pt}}\par}

% The two columns. \say is the whole unit: English left, French right.
%
% Fixed widths, not tabularx: tabularx scans its body for a literal
% \end{tabularx}, so wrapping it in a \newenvironment hides the terminator and
% the whole document ends up inside one unfinished cell. 84 + 5 + rule + 5 + 84
% is \textwidth on A4 with 15 mm margins.
\newcolumntype{S}{>{\setlength{\parskip}{1.1mm}}p{84mm}}
\newenvironment{say}{%
  \begin{tabular}{@{}S@{\hspace{5mm}}|@{\hspace{5mm}}S@{}}
}{%
  \end{tabular}\par}

% The deck's own shorthands. Frame titles are copied into this file VERBATIM by
% check_script.py, so anything main.tex can put in a title has to resolve here
% too -- \dB in "a $-5.1\dB$ discontinuity" is what first broke the build.
% Keep this list in step with the shorthands section of main.tex.
\newcommand{\dB}{\ensuremath{\,\mathrm{dB}}}  % \ensuremath: text here, math in the deck
\newcommand{\rr}{$R^2$}
\newcommand{\LA}{$L_{A,25\mathrm{s}}$}
\newcommand{\up}[1]{\textcolor{good}{\textbf{#1}}}
\newcommand{\down}[1]{\textcolor{bad}{\textbf{#1}}}
\newcommand{\wn}[1]{\textcolor{warn}{\textbf{#1}}}

% A line that must be said exactly, or must never be said.
\newcommand{\must}[1]{{\color{good}\textbf{#1}}}
\newcommand{\never}[1]{{\color{bad}\textbf{#1}}}
\newcommand{\care}[1]{{\color{warn}\textbf{#1}}}
\newcommand{\cue}[1]{\par\vspace{1mm}{\footnotesize\color{muted}\textit{#1}}\par}

% A boxed aside spanning the width: for questions to expect.
\newcommand{\aside}[2]{%
  \par\vspace{1.5mm}\noindent
  \colorbox{wash}{\parbox{\dimexpr\textwidth-2\fboxsep}{\vspace{1mm}
    \footnotesize\textbf{#1}\par\vspace{.7mm}#2\vspace{1mm}}}\par}


\scriptname{30 s a slide \; $\cdot$ \; EN spoken, FR pour comprendre}

\begin{document}
\selectlanguage{english}

\begin{center}
  {\LARGE\bfseries\color{accent} Speaker script --- short}\\[1.5mm]
  {\large Thirty seconds a slide}\\[3mm]
  {\color{muted}36 content slides \; $\cdot$ \; 8 section signposts \; $\cdot$ \;
   \textbf{18--19 minutes} \; $\cdot$ \; the long version is \texttt{script.pdf}}
\end{center}

\vspace{4mm}
\noindent
\colorbox{wash}{\parbox{\dimexpr\textwidth-2\fboxsep}{\vspace{1.5mm}
\textbf{How this differs from \texttt{script.pdf}.} Same talk, cut to fit. Each
English cell is about \textbf{sixty words} --- twenty-five seconds read aloud,
which leaves five a slide for the pause, the pointing and the breath. \must{Every caveat survived the cut.} What went is elaboration: the
second example, the reason behind the reason. A talk can be short or long; it
cannot be short and overclaiming, because \must{the sentence dropped under time
pressure is always the qualifying one}.

\medskip
\selectlanguage{french}
\textbf{Ce qui change par rapport à \texttt{script.pdf}.} Même exposé, coupé.
Chaque cellule anglaise fait environ \textbf{soixante mots} --- vingt-cinq
secondes à voix haute, ce qui laisse cinq secondes par diapo pour la pause, le
geste et la respiration. \must{Toutes les réserves ont survécu à
la coupe.} Ce qui saute, c'est le développement~: le deuxième exemple, la raison
derrière la raison. Un exposé peut être court ou long~; il ne peut pas être court
et exagéré, car \must{la phrase qu'on saute quand le temps manque est toujours
celle qui nuance}.
\selectlanguage{english}
\vspace{1.5mm}}}

\vspace{3mm}
\noindent
{\footnotesize\color{muted}Times are cumulative. \textbf{Section pages (3, 7, 10,
14, 21, 27, 37, 42) get one sentence each} --- name the section and move on. If
you are still over at slide 30, drop 32 and 33: they are engineering detail and
nothing later depends on them.}

\clearpage
% ============================================================================
\act{Act I --- the question, and what constrains it}{slides 1--13 \; $\cdot$ \; 0--5 min}

\slide{1}{Title}{0:00}
\begin{say}
Good morning. Our subject is what a smartphone protocol can, and cannot,
establish about urban noise.

One thing first. \must{This is not a noise map of Hanoi.} We cannot produce one.
It is a study of what this method actually supports.

Three months, 363 measurements, three ordinary phones --- and three negative
results that hold up better than the positive one.
&
Bonjour. Notre sujet~: ce qu'un protocole smartphone peut établir, et ne peut pas
établir, sur le bruit urbain.

Une chose d'abord. \must{Ce n'est pas une carte de bruit de Hanoï.} Nous sommes
incapables d'en produire une. C'est une étude de ce que cette méthode supporte
réellement.

Trois mois, 363 mesures, trois téléphones ordinaires --- et trois résultats
négatifs qui tiennent mieux que le positif.
\end{say}

\slide{2}{Outline}{0:30}
\begin{say}
Three parts. What we measured and what it supports. What our own audit found in
it. And where it goes --- an Android application, and a second city.

The audit is the part you may not expect. \must{It is us auditing ourselves.}
&
Trois parties. Ce que nous avons mesuré et ce que ça supporte. Ce que notre
propre audit y a trouvé. Et la suite --- une application Android, et une deuxième
ville.

L'audit est la partie inattendue. \must{C'est nous qui nous auditons.}
\end{say}
\cue{Section page 3 --- ``First, the context.'' Then move on.}

\slide{4}{The gap this project sits in}{1:05}
\begin{say}
Hanoi is dense, growing fast, and its traffic is motorcycles rather than cars.

Vietnam has no national noise-mapping programme, and exactly one instrumented
campaign has ever been published for Hanoi. It measured in 2005.

Class 1 meters and commercial propagation software are out of reach for most
local budgets. So: how far does a protocol built from three consumer phones and
open data actually get you?
&
Hanoï est dense, en croissance rapide, et son trafic est fait de deux-roues.

Le Vietnam n'a pas de programme national de cartographie du bruit, et une seule
campagne instrumentée a été publiée pour Hanoï. Elle mesurait en 2005.

Les sonomètres de classe 1 et les logiciels commerciaux sont hors de portée de la
plupart des budgets locaux. Donc~: jusqu'où va un protocole fait de trois
téléphones grand public et de données ouvertes~?
\end{say}

\slide{5}{The research question, and what is out of scope}{1:35}
\begin{say}
\must{``Can urban noise be attributed to the form of the city, and to the type of
vehicle passing through it?''} That is our question from June, word for word.
\must{We did not rewrite it after the answer came back.}

And the answer is no, on both halves. Morphology adds nothing the strict
protocols can see. Vehicle type: no --- and the chain meant to measure it does not
work.
&
\must{«~Le bruit urbain peut-il être attribué à la forme de la ville, et au type
de véhicule qui y circule~?~»} C'est notre question de juin, mot pour mot.
\must{Nous ne l'avons pas réécrite après avoir eu la réponse.}

Et la réponse est non, sur les deux moitiés. La morphologie n'apporte rien que
les protocoles stricts puissent voir. Le type de véhicule~: non --- et la chaîne
censée le mesurer ne fonctionne pas.
\end{say}

\slide{6}{Related work --- ten references, selected by what they are used for}{2:05}
\begin{say}
The bibliography is organised by what each reference is \emph{for}, not by date.

One paper deserves a sentence. Roberts and colleagues, 2017: spatial
cross-validation must exclude the area your features are built from.
\must{That is the paper that made us withdraw our own headline number.} It comes
back in part two.
&
La bibliographie est organisée par la \emph{fonction} de chaque référence, pas par
date.

Un article mérite une phrase. Roberts et collègues, 2017~: la validation croisée
spatiale doit exclure la zone d'où viennent vos variables. \must{C'est l'article
qui nous a fait retirer notre propre chiffre principal.} Il revient en partie
deux.
\end{say}
\cue{Section page 7 --- ``The data.''}

\slide{8}{The campaign}{2:40}
\begin{say}
363 points across three deliberately contrasted areas: a new peri-urban
development, the dense historic quarter, a transport corridor.

Left, the levels per site. Right, the hourly profile. \must{Those two contrasts
are what this protocol supports} --- the shape of the curves, not their height.

Three phones, twenty-five seconds a point, a ten-metre GPS gate. Both datasets
published.
&
363 points sur trois zones délibérément contrastées~: un développement
péri-urbain neuf, le quartier historique dense, un corridor de transport.

À gauche, les niveaux par site. À droite, le profil horaire. \must{Ces deux
contrastes sont ce que le protocole supporte} --- la forme des courbes, pas leur
hauteur.

Trois téléphones, vingt-cinq secondes par point, une porte GPS à dix mètres. Les
deux jeux publiés.
\end{say}

\slide{9}{Metrology --- the constraint everything else is built around}{3:10}
\begin{say}
This slide is the spine of everything that follows.

The three phones were calibrated against \emph{each other}. Never against a
reference. \must{A bias common to all three is invisible in our data by
construction.}

So: no absolute level as a fact, no compliance claim. Contrasts survive, because
a common bias cancels in a difference.

And the quantity is a twenty-five second integral. \must{Not $L_\mathrm{den}$,
which is a year.}
&
Cette diapo est la colonne vertébrale de tout ce qui suit.

Les trois téléphones ont été calibrés \emph{les uns sur les autres}. Jamais contre
une référence. \must{Un biais commun aux trois est invisible dans nos données par
construction.}

Donc~: aucun niveau absolu comme un fait, aucune revendication de conformité. Les
contrastes survivent, car un biais commun s'annule dans une différence.

Et la grandeur est une intégrale sur vingt-cinq secondes. \must{Pas
$L_\mathrm{den}$, qui est une année.}
\end{say}
\cue{Section page 10 --- ``How it was evaluated.'' \; \textbf{Expect the calibration
question here}: yes, half a day beside a class 2 meter would settle it; no, it was
not done. Slide 41 says so on the record.}

\slide{11}{The chain, end to end}{3:45}
\begin{say}
Left to right, once. Three sources: the survey, OpenStreetMap, the traffic
videos. Then cleaning, features, counting. Then evaluation. Then the outputs.

The line to remember is the last one. \must{The delivered model is chosen by
code, not by us.} Script 04 writes which model won; script 07 reads that flag.
Nothing downstream can quietly substitute another.
&
De gauche à droite, une fois. Trois sources~: l'enquête, OpenStreetMap, les vidéos
de trafic. Puis nettoyage, variables, comptage. Puis évaluation. Puis les sorties.

La ligne à retenir est la dernière. \must{Le modèle livré est choisi par le code,
pas par nous.} Le script 04 écrit quel modèle a gagné~; le 07 lit ce drapeau. Rien
en aval ne peut en substituer un autre.
\end{say}

\slide{12}{The evaluation protocol --- and why it is the whole story}{4:15}
\begin{say}
If you take one methodological slide, take this one.

Every model is scored under three splits, on identical folds. A permissive one.
The reference --- buffered leave-one-out at three hundred metres. And
leave-one-site-out.

Why three hundred? \must{Because the features aggregate over a disc of radius
three hundred metres.} Group closer points together and the model is tested on
near-twins of its training data. \must{The buffer must equal the feature radius.}
&
Si vous ne retenez qu'une diapo de méthode, retenez celle-ci.

Chaque modèle est évalué sous trois découpages, sur des plis identiques. Un
permissif. Le découpage de référence --- leave-one-out tamponné à trois cents
mètres. Et le leave-one-site-out.

Pourquoi trois cents~? \must{Parce que les variables agrègent sur un disque de
rayon trois cents mètres.} Regroupez des points plus proches et le modèle est
testé sur des quasi-jumeaux. \must{Le tampon doit égaler le rayon des variables.}
\end{say}

\slide{13}{Eight models, three baselines, one ablation}{4:45}
\begin{say}
Nothing is compared against nothing. The baselines are the point.

Mean per site and hour is a hard baseline, and we said so before running it.

And the hybrid --- physics plus machine learning --- is \must{the architecture we
ourselves had recommended}. We put it in the list and let it lose. That is the
whole point of the next section.
&
Rien n'est comparé à rien. Les baselines sont l'essentiel.

La moyenne par site et par heure est une baseline exigeante, et nous l'avons dit
avant de la lancer.

Et l'hybride --- physique plus apprentissage --- c'est \must{l'architecture que nous
avions nous-mêmes recommandée}. Nous l'avons mise dans la liste et laissée perdre.
C'est tout l'enjeu de la section suivante.
\end{say}
\cue{Section page 14 --- ``The results.''}

\clearpage
% ============================================================================
\act{Act II --- the results, and the audit}{slides 15--26 \; $\cdot$ \; 5--11 min}

\slide{15}{The ranking inverts --- and that is the result}{5:20}
\begin{say}
Give the room ten seconds before you speak.

\care{Read across, not down.} On the permissive split the elaborate hybrid wins.
On the two strict splits it collapses. The green line, three parameters, barely
moves.

\must{The ranking reverses between a permissive split and one that tests
generalisation.} That is the result --- not which model won, but that the answer
depends on the split.
&
Laissez la salle dix secondes avant de parler.

\care{Lisez en travers, pas de haut en bas.} Sur le découpage permissif, l'hybride
sophistiqué gagne. Sur les deux stricts, il s'effondre. La ligne verte, trois
paramètres, bouge à peine.

\must{Le classement s'inverse entre un découpage permissif et un découpage qui
teste la généralisation.} C'est ça le résultat --- pas quel modèle gagne, mais que
la réponse dépende du découpage.
\end{say}

\slide{16}{The delivered model}{5:50}
\begin{say}
What won is a line-source attenuation law. Two distances, three constants,
\must{no morphology at all}. Energy falls as one over $d$ --- a street is a line,
not a point.

One honesty note: the background term fits to a boundary. \must{It is not
identifiable from this data.}

And say the error out loud. \care{A typical prediction is wrong by five
decibels.} Enough for a contrast, nowhere near an assessment.
&
Ce qui gagne est une loi d'atténuation de source linéique. Deux distances, trois
constantes, \must{aucune morphologie}. L'énergie décroît en un sur $d$ --- une rue
est une ligne, pas un point.

Une note d'honnêteté~: le terme de fond converge vers une borne. \must{Il n'est
pas identifiable avec ces données.}

Et dites l'erreur à voix haute. \care{Une prédiction typique se trompe de cinq
décibels.} Assez pour un contraste, très loin d'une évaluation.
\end{say}

\slide{17}{``$R^2 = 0.25$ --- isn't that very low?''}{6:20}
\begin{say}
This question always comes. Answer it first.

Two points a few metres apart at the same hour must get the same prediction from
any spatial model. \must{Whatever they disagree by is variance no model can
explain.} They differ by four decibels.

At the forty-metre scale the map works at, \must{the ceiling is about 0.54}. We
reach 0.246 --- half of what is attainable, not a tenth.
&
Cette question vient toujours. Répondez-y en premier.

Deux points à quelques mètres, à la même heure, doivent recevoir la même
prédiction. \must{Leur désaccord est une variance qu'aucun modèle ne peut
expliquer.} Ils diffèrent de quatre décibels.

À l'échelle de quarante mètres de la carte, \must{le plafond est d'environ 0,54}.
Nous atteignons 0,246 --- la moitié de l'atteignable, pas le dixième.
\end{say}
\cue{If you have the time, add the caveat: same-hour pairs may be on different
days, so this overestimates measurement noise. If you do not, it is on the slide.}

\slide{18}{The map --- and the caveat that goes with it}{6:50}
\begin{say}
The predicted grid over the three measured areas. Forty-metre cells.

Two caveats, and I want to be first with both. \care{The grid step is not the
accuracy} --- forty-metre cells from a model with five decibels of error.

And \care{there is no hour}: all seventeen hourly columns are identical. The
kernel has no hour term.

The colours are QCVN bands, descriptive. \never{Not a compliance statement.}
&
La grille prédite sur les trois zones mesurées. Cellules de quarante mètres.

Deux réserves, et je veux être le premier à les dire. \care{Le pas de la grille
n'est pas la précision} --- quarante mètres, pour un modèle à cinq décibels
d'erreur.

Et \care{il n'y a pas d'heure}~: les dix-sept colonnes horaires sont identiques.
Le noyau n'a pas de terme horaire.

Les couleurs sont les bandes QCVN, descriptives. \never{Pas une conformité.}
\end{say}

\slide{19}{What the learned model was leaning on}{7:20}
\begin{say}
Two thirds of the learned model's gain goes to morphology.

And it buys nothing: morphology alone scores 0.041 under the reference protocol.

\must{Where a model spends its capacity and where it earns its score are
different questions.} Only the second is a result.
&
Deux tiers du gain du modèle appris vont à la morphologie.

Et ça n'achète rien~: la morphologie seule obtient 0,041 sous le protocole de
référence.

\must{Où un modèle dépense sa capacité et où il gagne son score sont deux
questions différentes.} Seule la seconde est un résultat.
\end{say}

\slide{20}{The three negative results}{7:50}
\begin{say}
These three are the contribution. Say it in those words.

One: three physical parameters beat every learned model. Capacity does not pay.

Two: cross-city transfer fails. Not a calibration offset --- the Uganda model
ranks quiet and loud \must{the wrong way round}. Against which one local distance
term reaches 0.240.

Three: vehicle flow does not explain the levels --- and our own audit says that is
an instrument result.
&
Ces trois-là sont la contribution. Dites-le avec ces mots.

Un~: trois paramètres physiques battent tous les modèles appris. La capacité ne
paie pas.

Deux~: le transfert entre villes échoue. Pas un décalage de calibration --- le
modèle ougandais classe le calme et le bruyant \must{à l'envers}. Face à quoi un
seul terme de distance local atteint 0,240.

Trois~: le flux de véhicules n'explique pas les niveaux --- et notre audit dit que
c'est un résultat instrumental.
\end{say}
\cue{Section page 21 --- ``Now we audit ourselves.''}

\slide{22}{We found one of our own published results was wrong}{8:25}
\begin{say}
In August we withdrew our own headline number.

A model scored 0.45 under what we then called honest spatial cross-validation. A
simulation was built on it. The grouping was on hundred-and-ten-metre cells while
the features span three hundred. \must{It leaked.}

What replaced it is 0.246. The retracted deck and map are archived \must{with
their reason}, not deleted.
&
En août nous avons retiré notre propre chiffre principal.

Un modèle obtenait 0,45 sous ce que nous appelions alors une validation croisée
spatiale honnête. Une simulation avait été construite dessus. Le regroupement se
faisait sur des cellules de cent dix mètres, les variables sur trois cents.
\must{Il y avait fuite.}

Ce qui l'a remplacé, c'est 0,246. La présentation et la carte retirées sont
archivées \must{avec leur raison}, pas supprimées.
\end{say}

\slide{23}{A second audit, on the data itself}{8:55}
\begin{say}
The first audit checked the method. This one checks the measurements.

The top half matters as much as the bottom: 363 raw rows to 363 published, nothing
dropped. Every README figure matches the metrics file exactly.

\must{That is what makes the bottom half credible rather than alarming.} Three
findings, none of them known when the results were written.
&
Le premier audit vérifiait la méthode. Celui-ci vérifie les mesures.

La moitié haute compte autant que la basse~: 363 lignes brutes pour 363 publiées,
rien de perdu. Chaque chiffre du README correspond exactement au fichier de
métriques.

\must{C'est ça qui rend la moitié basse crédible plutôt qu'alarmante.} Trois
constats, aucun connu quand les résultats ont été écrits.
\end{say}

\slide{24}{Finding 1 --- a $-5.1\dB$ discontinuity on 27 June}{9:25}
\begin{say}
On 27 June the level distribution steps down by five decibels, and it survives
controlling for site, hour and class.

Be precise. \never{This does not show the $R^2$ is inflated.} It shows a large,
undocumented discontinuity the model has no term for.

The worrying part is that it is \care{spatially structured} --- so if it is
instrumental, a model can learn it as morphology. Status: open.
&
Le 27 juin, la distribution descend d'une marche de cinq décibels, et ça survit au
contrôle par site, heure et classe.

Soyez précis. \never{Ça ne montre pas que le $R^2$ est gonflé.} Ça montre une
discontinuité large et non documentée que le modèle n'a aucun terme pour
représenter.

L'inquiétant est qu'elle soit \care{spatialement structurée} --- si elle est
instrumentale, un modèle peut l'apprendre comme de la morphologie. Statut~:
ouvert.
\end{say}

\slide{25}{Findings 2 and 3 --- the video chain, and the intervals}{9:55}
\begin{say}
The detector reports fifty-seven percent cars in Hanoi. Reality is the inverse.
And video-to-measurement gaps reach a hundred and sixty seconds, for a
twenty-five second sample.

So negative result three is far more likely an \care{instrument failure} than
acoustics.

Finding three is presentational: the headline ranking is \must{not statistically
separated}. The intervals overlap almost entirely.
&
Le détecteur annonce cinquante-sept pour cent de voitures à Hanoï. La réalité est
l'inverse. Et les écarts vidéo--mesure atteignent cent soixante secondes, pour un
échantillon de vingt-cinq.

Donc le résultat négatif trois est bien plus vraisemblablement un \care{échec
instrumental} que de l'acoustique.

Le constat trois est de présentation~: le classement principal n'est \must{pas
séparé statistiquement}. Les intervalles se recouvrent presque entièrement.
\end{say}

\slide{26}{What the audit leaves standing}{10:25}
\begin{say}
The protocol holds. The ranking inversion holds. Negative results one and two
hold. Provenance is lossless and checkable.

What needs work: diagnose the June discontinuity, reframe result three, carry the
intervals everywhere.

And here is the hinge. \must{Everything the audit could not repair comes back to
one thing: we did not control the instrument.}
&
Le protocole tient. L'inversion tient. Les résultats négatifs un et deux tiennent.
La provenance est sans perte et vérifiable.

Ce qui reste~: diagnostiquer la discontinuité de juin, requalifier le résultat
trois, propager les intervalles partout.

Et voici la charnière. \must{Tout ce que l'audit n'a pas pu réparer revient à une
seule chose~: nous ne contrôlions pas l'instrument.}
\end{say}
\cue{Section page 27 --- ``So we built one.''}

\clearpage
% ============================================================================
\act{Act III --- the instrument}{slides 28--36 \; $\cdot$ \; 11--15 min}

\slide{28}{Three months, in one picture}{11:00}
\begin{say}
Do not read the boxes. Point at the two raised labels --- the two retractions ---
then at the last one.

\must{Every phase ended by rejecting its own preferred answer.} The transfer we
set out to do. The $R^2$ we had presented. The architecture we had recommended
ourselves.

Each rejection pointed at the same root cause. So the last three weeks built an
instrument.
&
Ne lisez pas les cases. Pointez les deux étiquettes surélevées --- les deux
retraits --- puis la dernière.

\must{Chaque phase s'est terminée en rejetant sa propre réponse préférée.} Le
transfert visé. Le $R^2$ présenté. L'architecture que nous recommandions.

Chaque rejet pointait la même cause racine. Donc les trois dernières semaines ont
construit un instrument.
\end{say}

\slide{29}{Why build an application at all}{11:30}
\begin{say}
Four things the campaign could not do.

Control the instrument --- the app we used is closed. Measure and record in one
session --- two apps meant the level of one street and the audio of another. Scale
beyond three people. And record what the device even was.

Point two is a real platform constraint: Android will not give the microphone to
two consumers at once.
&
Quatre choses que la campagne ne pouvait pas faire.

Contrôler l'instrument --- l'application utilisée est fermée. Mesurer et
enregistrer en une session --- deux applications, donc le niveau d'une rue et
l'audio d'une autre. Passer l'échelle au-delà de trois personnes. Et enregistrer
ce qu'était l'appareil.

Le point deux est une vraie contrainte~: Android ne donne pas le microphone à deux
consommateurs à la fois.
\end{say}

\slide{30}{The application}{12:00}
\begin{say}
One sentence each. Collect --- both forms, offline. Measure --- twenty-five
seconds, A-weighted. Map --- the predicted level where you stand. Results --- read
from the metrics file, negative results included.

\must{And say this plainly: these are drawings from the source code, not
screenshots.} The app has never run in front of a sound source. That is the
limitation four slides from now.
&
Une phrase chacun. Collecter --- les deux formulaires, hors ligne. Mesurer ---
vingt-cinq secondes, pondéré A. Carte --- le niveau prédit là où vous êtes.
Résultats --- lus depuis le fichier de métriques, négatifs compris.

\must{Et dites-le franchement~: ce sont des dessins faits depuis le code, pas des
captures.} L'application n'a jamais tourné devant une source sonore. C'est la
limite de la diapo dans quatre.
\end{say}

\slide{31}{What was built --- phase by phase}{12:30}
\begin{say}
Summarise with a supervisor; slow down with the people who inherit it. \care{This
is the slide they should photograph.}

Phases one and two were deliberately not separated, for the microphone reason.

Phase four got simpler on contact with reality: the published grid is already the
kernel's output, so reading the nearest cell \emph{is} the prediction.

\must{Phase six is not built, and I will not hide it.}
&
Résumez face à un encadrant~; ralentissez face à ceux qui en héritent. \care{C'est
la diapo qu'ils devraient photographier.}

Les phases un et deux n'ont délibérément pas été séparées, pour le microphone.

La phase quatre s'est simplifiée au contact du réel~: la grille publiée est déjà
la sortie du noyau, donc lire la cellule la plus proche \emph{est} la prédiction.

\must{La phase six n'est pas construite, et je ne le cache pas.}
\end{say}

\slide{32}{Architecture --- and the one design rule that matters}{13:00}
\begin{say}
One rule. \must{The application does not carry a copy of the study.} The build
copies the map, the measurements and the metrics file out of the repository. Re-run
the pipeline, rebuild, and the app moves with it.

Same principle, third time: the report reads the metrics file, the dashboard reads
it, now the app reads it. \must{Every published number has exactly one home.}
&
Une règle. \must{L'application ne porte pas une copie de l'étude.} La compilation
copie la carte, les mesures et le fichier de métriques depuis le dépôt. Relancez
le pipeline, recompilez, l'application suit.

Même principe, troisième fois~: le rapport lit le fichier de métriques, le tableau
de bord le lit, maintenant l'application le lit. \must{Chaque chiffre publié a
exactement un domicile.}
\end{say}
\cue{\textbf{Behind time?} Cut this slide and the next. Nothing later depends on
either.}

\slide{33}{Three engineering results worth reporting}{13:30}
\begin{say}
Three things that each cost a day and are invisible in a results table.

The middle one is worth dwelling on. \must{The same mistake was made twice.}
Scaling total energy by a traffic multiplier double-counts the background. We
corrected it in July, then rebuilt it in August knowing about it.

\must{Knowing a mistake is not being immune to it.} A test caught it.
&
Trois choses qui ont coûté un jour chacune et sont invisibles dans un tableau.

Celle du milieu mérite qu'on s'y arrête. \must{La même erreur a été commise deux
fois.} Multiplier l'énergie totale par un facteur de trafic compte le fond deux
fois. Corrigée en juillet, reconstruite en août en la connaissant.

\must{Connaître une erreur ne rend pas immunisé.} Un test l'a rattrapée.
\end{say}

\slide{34}{What the app measures --- and what it refuses to claim}{14:00}
\begin{say}
The app's level is a \emph{different quantity} from the campaign's. It goes in
its own field. \must{The two never share a column.}

Verified end to end against the live server: a fix, a measurement, a submission
accepted.

And the important one, unverified. \must{The level has never been checked against
a sound source} --- which, to be even-handed, \must{is equally true of the
campaign's numbers.}
&
Le niveau de l'application est une \emph{grandeur différente} de celle de la
campagne. Il va dans son propre champ. \must{Les deux ne partagent jamais une
colonne.}

Vérifié de bout en bout contre le serveur réel~: un point GPS, une mesure, une
soumission acceptée.

Et l'important, non vérifié. \must{Le niveau n'a jamais été confronté à une source
sonore} --- ce qui, pour être équitable, \must{vaut aussi pour les chiffres de la
campagne.}
\end{say}

\slide{35}{Driving GAMA from the phone --- how it is wired}{14:30}
\begin{say}
GAMA has no Android port, but it exposes a WebSocket API. You start it headless
with a socket, and the app is a client. \must{The phone is a remote control and a
screen. It runs nothing.}

It redraws the model's own display, layer by layer in declaration order.

What it buys is the half a precomputed grid cannot hold: mitigation scenarios,
construction, the fleet by hour.
&
GAMA n'a pas de portage Android, mais expose une API WebSocket. On la lance en
headless avec un socket, et l'application est un client. \must{Le téléphone est
une télécommande et un écran. Il n'exécute rien.}

Elle redessine l'affichage propre du modèle, couche par couche dans l'ordre de
déclaration.

Ce que ça apporte, c'est la moitié qu'une grille précalculée ne peut pas porter~:
mitigation, chantiers, flotte par heure.
\end{say}

\slide{36}{Before this app goes to strangers --- three blockers}{15:00}
\begin{say}
Three things stop this being handed out, and \must{none is a bug.}

Whoever owns the Kobo account holds strangers' positions and answers for them. It
must be institutional. \must{A decision, not a code change.}

A consent screen is not an ethics review, and a public campaign is a new campaign.

And the GAMA server needs hosting, with TLS.
&
Trois choses empêchent de la distribuer, et \must{aucune n'est un bug.}

Celui qui détient le compte Kobo détient les positions d'inconnus et en répond. Il
doit être institutionnel. \must{Une décision, pas une modification.}

Un écran de consentement n'est pas un avis éthique, et une campagne publique est
une nouvelle campagne.

Et le serveur GAMA demande de l'hébergement, avec TLS.
\end{say}
\cue{Section page 37 --- ``And next.''}

\clearpage
% ============================================================================
\act{Act IV --- what comes next}{slides 38--44 \; $\cdot$ \; 15--19 min}

\slide{38}{Clermont-Ferrand --- what it does and does not change}{15:35}
\begin{say}
Correct one expectation before anyone forms it. \must{There will be no sound level
meter at Clermont either.} The calibration gap does not close.

So the honest description is \must{maintenance and deployment}: keep the
application alive, then get it collecting in a second city.

That buys the strongest experiment available. \must{The same binary in both
cities} --- the instrument stops being a confound.
&
Corrigez une attente avant qu'elle se forme. \must{Il n'y aura pas non plus de
sonomètre à Clermont.} L'écart de calibration ne se referme pas.

Donc la description honnête, c'est \must{maintenance et déploiement}~: garder
l'application vivante, puis la faire collecter dans une deuxième ville.

Ça achète l'expérience la plus forte disponible. \must{Le même binaire dans les
deux villes} --- l'instrument cesse d'être un facteur de confusion.
\end{say}

\slide{39}{The comparative deployment --- and why it is the decisive test}{16:05}
\begin{say}
Negative result two says transfer fails. But there, \care{the instrument, the
protocol and the feature conventions all differed too.} It is an over-determined
result.

One application in both cities, and \must{a transfer failure means what it says.}

And Clermont falls under EU strategic mapping --- \must{the first official
instrumented reference map} we can compare against.
&
Le résultat négatif deux dit que le transfert échoue. Mais là, \care{l'instrument,
le protocole et les conventions différaient aussi.} Résultat sur-déterminé.

Une seule application dans les deux villes, et \must{un échec de transfert veut
dire ce qu'il dit.}

Et Clermont relève de la cartographie stratégique européenne --- \must{la première
carte de référence officielle instrumentée} à laquelle nous comparer.
\end{say}

\slide{40}{The programme, staged}{16:35}
\begin{say}
Five stages, sequenced by dependency, not by calendar. Every stage has a gate, and
\must{the gates are falsifiable} --- a roadmap without exit conditions is a wish
list.

Point at the physics gate: it is written so a negative answer is still a result.

And be honest that stage four depends on things nobody has confirmed.
\never{Do not present it as scheduled.}
&
Cinq étapes, séquencées par dépendance, pas par calendrier. Chaque étape a une
porte, et \must{les portes sont falsifiables} --- une feuille de route sans
conditions de sortie est une liste de vœux.

Pointez la porte de la physique~: elle est écrite pour qu'une réponse négative
reste un résultat.

Et soyez honnête~: l'étape quatre dépend de choses non confirmées. \never{Ne la
présentez pas comme planifiée.}
\end{say}

\slide{41}{The two things a reviewer will press on}{17:05}
\begin{say}
Two objections, better raised by us than by a reviewer.

No absolute reference. A reviewer will not ask whether we \emph{knew} --- they will
ask whether one reading beside a class 2 meter was unobtainable. \must{A half-day
rental would settle it.}

And three sites are three typologies. \must{The hard limit is the design, not the
$R^2$.}
&
Deux objections, mieux soulevées par nous que par un relecteur.

Pas de référence absolue. Un relecteur ne demandera pas si nous \emph{savions} ---
il demandera si une mesure à côté d'un sonomètre de classe 2 était inatteignable.
\must{Une demi-journée de location y répondrait.}

Et trois sites, c'est trois typologies. \must{La limite dure, c'est le plan
d'échantillonnage, pas le $R^2$.}
\end{say}
\cue{Section page 42 --- ``To close.''}

\slide{43}{What this work established}{17:40}
\begin{say}
Four findings, and six practices that make them trustworthy.

If the room is technical, close on the right column. Those six transfer to any
spatial machine learning project.

The buffer matches the feature radius. The delivered model is chosen by code.
Every number has one home. Retracted work is archived with its reason.

\must{The deliverable is a protocol, an instrument and an honest error bar. Not a
noise map of Hanoi.}
&
Quatre constats, et six pratiques qui les rendent fiables.

Si la salle est technique, terminez sur la colonne de droite. Ces six-là se
transfèrent à tout projet de machine learning spatial.

Le tampon égale le rayon des variables. Le modèle livré est choisi par le code.
Chaque chiffre a un domicile. Le travail retiré est archivé avec sa raison.

\must{Le livrable, c'est un protocole, un instrument et une barre d'erreur
honnête. Pas une carte de bruit de Hanoï.}
\end{say}
\cue{Say the last line slowly.}

\slide{44}{Thank you}{18:10}
\begin{say}
Thank you. Everything on these slides rebuilds from the repository --- four make
commands, from two published CSV files.

I would rather offer you that than a summary.
&
Merci. Tout se reconstruit depuis le dépôt --- quatre commandes make, à partir de
deux fichiers CSV publiés.

Je préfère vous offrir ça qu'un résumé.
\end{say}

\vspace{6mm}
\noindent
\colorbox{wash}{\parbox{\dimexpr\textwidth-2\fboxsep}{\vspace{1.5mm}
\textbf{Questions.} Four backup slides: \textbf{45} every formula, \textbf{46} the
forest plot of intervals, \textbf{47} the bias interval, \textbf{48} what is
published. The four answers to have ready are in \texttt{script.pdf}, and they are
short enough to memorise: \must{``It is not calibrated.''} ---
\must{``We did not test that.''} --- \must{``That is descriptive, not a compliance
claim.''} --- \must{``It is in the repository, with its reason.''}

\medskip
\selectlanguage{french}
\textbf{Questions.} Quatre diapos de secours~: \textbf{45} les formules,
\textbf{46} les intervalles, \textbf{47} le biais, \textbf{48} ce qui est publié.
Les quatre réponses à avoir prêtes sont dans \texttt{script.pdf}, assez courtes
pour être apprises.
\selectlanguage{english}
\vspace{1.5mm}}}

\end{document}
